\documentclass{article}
\usepackage{graphicx} % Required for inserting images
\usepackage{amsfonts, amsmath} % For \mathbb

\newtheorem{theorem}{Theorem}
\newtheorem{exercise}{Exercise}
\newcommand{\trace}{\operatorname{tr}}
\newcommand{\ev}{\operatorname{ev}}
\newcommand{\coev}{\operatorname{coev}}
\newcommand{\im}{\operatorname{im}}
\newcommand{\End}{\operatorname{End}}



\title{Coordinate-free Mathematics}
\author{Paul Gustafson}
\date{May 2025}

\begin{document}


\maketitle
\section{Linear Algebra}

Linear algebra studies linear maps between vector spaces. A \textbf{vector space} $V$ over a field $k$ is a set equipped with vector addition and scalar multiplication operations satisfying straightforward compatibility conditions. A \textbf{linear map} $f: V \to U$ between $k$-vector spaces is a map of underlying sets such that $f(\alpha v + \beta w) = \alpha f(v) + \beta f(w)$ for all $\alpha, \beta \in k$ and $v, w \in V$. Examples include differential operators, group representations, coefficient matrices of linear systems of equations, and local approximations of globally nonlinear functions.

The \textbf{free $k$-vector space on a set $X$}, denoted $kX$, is the set of finitely supported set-theoretical functions $X \to k$ with pointwise addition and scalar multiplication. The \textbf{direct sum} vector space $V \oplus W$ is the set $V \times W$ with pointwise addition and scalar multiplication defined by $\alpha (v \oplus w) = \alpha v \oplus \alpha w$. The \textbf{quotient} of $V$ by a subspace $W$, denoted $V / W$, is the set of equivalence classes of $V$ modulo the relation $x \sim y$ iff $x - y \in W$. The \textbf{tensor product} $V \otimes W$ is the vector space $k(V \times W) / N$ where $N$ is the subspace generated by the obvious linearity relations on $V$ and $W$.

\begin{exercise} An \textbf{isomorphism} of vector spaces $V \cong W$ consists of a pair of linear maps $f : V \to W$ and $g : W \to V$ such that $f \circ g = 1_W$ and $g \circ f = 1_V$. Prove the following:
\begin{itemize}
\item $k(X \sqcup Y) \cong kX \oplus kY$
\item $k(X \times Y) \cong kX \otimes kY$
\item The first isomorphism theorem for vector spaces (and second and third if you want).
\item Find universal properties determining up to canonical isomorphism the free $k$-vector space, direct sum, quotient, tensor product, and dual space (defined below), respectively.
\end{itemize}
\end{exercise}

The \textbf{dual space} $V^*$ is the space of linear maps from a vector space $V$ to its ground field $k$. The evaluation map $\ev_V : V^* \otimes V \to k$ is the linear extension of the evaluation morphism $f \otimes x \mapsto f(x)$. We say that $V$ is \textbf{finite dimensional} if there is a co-evaluation morphism $\coev_V: k \to V \otimes V^*$ satisfying the snake equations $(1_{V} \otimes \ev_V) \circ (\coev_V \otimes 1_V) = 1_V$
and
$(\ev_V \otimes 1_{V^*}) \circ (1_{V^*} \otimes \coev_{V}) = 1_{V^*}$. We define the \textbf{dimension} $\dim(V) = \trace(1_V)$, where $\trace(f) : \End(V) \to k$ is given by $f \mapsto (\ev_V \circ s_{V, V^*} \circ (f \otimes 1_{V^*}) \circ \coev_V)(1)$ where $s_{X, Y} : X \otimes Y \to Y \otimes X$ is the symmetric braiding $x \otimes y \mapsto y \otimes x$.



\begin{exercise} Prove the following for finite-dimensional vector spaces $V, W$ and linear maps $f, g : V \to V$.
    \begin{itemize}
        \item $\trace(f g) = \trace(g f)$
        \item $\trace(f \oplus g) = \trace(f) + \trace(g)$
        \item $\trace(f \otimes g) = \trace(f) \trace(g)$
        \item $\dim(V) \in \mathbf{N}$
        \item $V \cong W$ iff $\dim(V) = \dim(W)$
    \end{itemize}
\end{exercise}


The simplest linear maps are scalar multiples of the identity. Although most linear maps are not of this form, it is useful to decompose them into sums of scalar maps when possible. One condition under which we can decompose a linear map into a direct sum of scalars is if the map is self-adjoint. Suppose $k = \mathbb{C}$, the field of complex numbers. A Hermitian inner product $\langle \cdot, \cdot \rangle: V \times V \to \mathbb{C}$ is a positive-definite sesquilinear form (linear in the first variable, anti-linear in the second variable, and conjugate symmetric). A linear map $f : V \to V$ is \textbf{self-adjoint} if $\langle f v, w \rangle = \langle v, f w \rangle$ for all $v, w \in V$.

\begin{exercise}
Prove that any self-adjoint linear map $f$ on a finite-dimensional complex vector space $V$ is diagonalizable:
$$f = \sum \lambda_i p_i,$$
where $\sum p_i = 1_V$ and $p_i p_j = \delta_{ij} p_i$. Use this fact to prove the existence of the singular value decomposition for an arbitrary linear map between complex finite-dimensional vector spaces.
\end{exercise}

The scalars $\lambda_i$ above are called the eigenvalues of $f$, and the images of $p_i$ are the corresponding eigenspaces. More generally, a scalar $\lambda \in k$ is an \textbf{eigenvalue} for a linear map $f : V \to V$ if there exists a nonzero vector $v \in V$ for which $fv = \lambda v$. The \textbf{eigenspace} of $\lambda$ is
\[
\ker(f-\lambda 1_V),
\]
whose nonzero elements are eigenvectors; its dimension is called the \textbf{geometric multiplicity} of $\lambda$.


We can solve for the eigenvalues of a linear operator on a finite-dimensional vector space via the determinant. Define the tensor algebra and \textbf{exterior algebra of $V$} by
\[
T(V)=\bigoplus_{k\geq 0}V^{\otimes k},
\qquad
\Lambda V=T(V)/\langle v\otimes v:v\in V\rangle.
\]
The space $\Lambda^kV$ is the degree-$k$ component of $\Lambda V$.

\begin{exercise}
    Show that the highest graded component of $\Lambda V$ is one-dimensional.
\end{exercise}

The \textbf{determinant} $\det(f)$ is the coefficient of the linear map on the highest-graded component of $\Lambda V$ given by $v_1 \wedge \cdots \wedge v_{\dim(V)} \mapsto f(v_1) \wedge \cdots f(v_{\dim(V)})$.

\begin{exercise} Prove the following:
\begin{itemize}
    \item $\det(f \oplus g) = \det(f) \det(g)$
    \item If $V = \mathbb{R}^n$, then $\det(f)$ is the signed volume of $f(I^n)$ where $I = [0,1]$.
    \item Cramer's rule
\end{itemize}
\end{exercise}

For $f\in\End(V)$, the \textbf{characteristic polynomial} is the polynomial $\lambda \mapsto \det(f - \lambda 1_V)$.

\begin{exercise}
    Assume that the characteristic polynomial of $f$ splits over $k$. Prove the following for a linear operator $f$ on a finite-dimensional vector space $V$:
    \begin{itemize}

        \item Every eigenvalue of $f$ is a root of its characteristic polynomial.
        \item The geometric multiplicity of each eigenvalue is at most equal to its \textbf{algebraic multiplicity} (i.e. multiplicity as a root of the characteristic polynomial)
        \item $\trace(f) = \sum_i \lambda_i$ and $\det(f) = \prod_i \lambda_i$, where each eigenvalue $\lambda_i$ is repeated according to its algebraic multiplicity
        \item The operator $f$ satisfies its own characteristic polynomial (replacing $\lambda$ with $f$)
        \item The existence of the Jordan normal form for a linear operator on a finite-dimensional complex vector space
    \end{itemize}
\end{exercise}

\section{Calculus}
A \textbf{smooth manifold} is a topological space that is locally diffeomorphic to Euclidean space $\mathbb{R}^n$. A \textbf{smooth vector bundle} $p : E \to M$ consists of a smooth surjection $p$ between smooth manifolds $E$ (the total space) and $M$ (the base space) such that all fibers are isomorphic to a single vector space $V$, and smooth local trivializations exist (locally, the fiber bundle looks like a product space). A prototypical example is the \textbf{tangent bundle} $TM$ which consists of all tangent spaces of points of a given manifold. A \textbf{vector field} $X \in \Gamma(TM)$ is a smooth section of the tangent bundle (i.e. a smooth map $s: M \to TM$ such that $p \circ s = 1_M$).

\begin{exercise}
Write down precise definitions for the following:
\begin{itemize}
\item smooth manifold, smooth map, vector bundle, tangent bundle.
\item bundle map, direct sum, tensor product, and dual bundle
\item derivative of a smooth function between manifolds
\end{itemize}
\end{exercise}

To define integration in a coordinate-free way, we will integrate differential forms over chains. A \textbf{differential $n$-form} is a section of the $n$-th graded component of the exterior bundle $\Lambda T^*M$. A \textbf{singular $n$-simplex} is a smooth map $\sigma:\Delta^n \to M$ where $\Delta^n$ is the standard $n$-simplex. The collection of singular $n$-simplices in $M$ freely generates the abelian group $C_n(M)$ of \textbf{singular $n$-chains} in $M$. The \textbf{boundary map} $\partial : C_n(M) \to C_{n-1}(M)$ is defined on each simplex $\sigma:\Delta^n \to M$ by
$$\partial\sigma=\sum_{j=0}^{n}(-1)^j\sigma\circ f_j,$$ where $f_j : \Delta^{n-1} \to \Delta^n$ is the inclusion of the face opposite the $j$-th vertex of $\Delta^n$.

A \textbf{chain complex} $(A_n, d_n)_n$ is a sequence of abelian groups $A_n$ and homomorphisms $d_n : A_n \to A_{n-1}$ such that $d_{n-1} \circ d_n = 0$ for all $n$. The $n$-th \textbf{homology group} $H_n$ of the chain complex is
$$H_n = \ker(d_{n}) / \im(d_{n+1}).$$
The prototypical example is singular homology which corresponds to the chain complex of singular chains with the boundary map defined above.

Cohomology and cochains are defined similarly, using superscripts instead of subscripts and boundary maps $d^n : A^n \to A^{n+1}$ increase degree instead of decrease it. Singular cohomology corresponds to the cochain complex given by dualizing the singular chain complex above.

The \textbf{de Rham cohomology} of $M$ is
\[
H_{\mathrm{dR}}^n(M)
=
\ker(d\colon\Omega^n(M)\to\Omega^{n+1}(M))/
\operatorname{im}(d\colon\Omega^{n-1}(M)\to\Omega^n(M)).
\]

\begin{exercise}
Prove the following:
\begin{itemize}
\item Define the \textbf{exterior derivative} of an $n$-form, and show that differential forms on a manifold $M$ form a cochain complex.
\item Define the integral of an $n$-form on an $n$-chain. Prove that $\int_c d \omega = \int_{\partial c} \omega$ for any $n$-chain $c$ and differential $n-1$-form $\omega$.

\item Singular cohomology with real coefficients and de Rham cohomology are naturally isomorphic.

\item Bounded chain complexes of finite-dimensional vector spaces form a symmetric monoidal abelian category with duals, the trace of an endomorphism is its Lefschetz number, and the dimension of a chain complex is its Euler characteristic.
\end{itemize}
\end{exercise}

An \textbf{affine connection} is a map $\nabla : \Gamma(TM) \times \Gamma(TM) \to \Gamma(TM)$ such that $\nabla_{fX} Y = f \nabla_X Y$ and
$\nabla_X(fY) = (\partial_X f) Y + f \nabla_X Y$ for any scalar field $f$ on $M$.







\begin{thebibliography}{9}

\bibitem{halmos}
Paul R. Halmos,
\textit{Finite-Dimensional Vector Spaces},
2nd ed., Springer, 1974.

\bibitem{hoffman-kunze}
Kenneth Hoffman and Ray Kunze,
\textit{Linear Algebra},
2nd ed., Prentice Hall, 1971.

\bibitem{spivak}
Michael Spivak,
\textit{Calculus on Manifolds: A Modern Approach to Classical Theorems of Advanced Calculus},
Addison-Wesley, 1965.

\bibitem{lang}
Serge Lang,
\textit{Algebra},
3rd ed., Springer, 2002.


\bibitem{ponto-shulman}
Kate Ponto and Michael Shulman,
\textit{Traces in symmetric monoidal categories},
Expositiones Mathematicae, 32(3):248--273, 2014.

\bibitem{bott-tu}
Raoul Bott and Loring W. Tu,
\textit{Differential Forms in Algebraic Topology},
Graduate Texts in Mathematics, Vol. 82, Springer, 1982.

\end{thebibliography}


\end{document}
